\scanpage{10}
\mspage{24}
\margin{$/^{\circ}[\,]$ inutile, \unsure{résulte du} \deleted{\unsure{fait}} que $IB$ (\unsure{resp.} $IC$) top. nilpotent.}

\uline{contenant $IB$ resp. $IC$}\inserted{]}\unsure{,} enfin on suppose $IB$ et $IC$ fermés (condition automatiquement vérifiée si $B,C$ sont noethériens). \unsure{Soient}\deleted{\unsure{u}} $B_0=B/IB$,\deleted{et} $C_0=C/IC$, $u:\inserted{B\to C}$ un homomorphisme continu de $A$-algèbres, $u_i:B_i\to C_i$ l'homomorphisme \deleted{induit} déduit par passage au quotient (j'ai oublié d'introduire un système fondamental d'idéaux ouverts $I_i\subset I_0\inserted{=I}$ dans $A$, et les $B_i$ et $C_i$ en conséquence\inserted{enfin ce sont tous les $I_iB$ et $I_iC$ qui doivent être fermés}). Conditions équivalentes:

(i\inserted{bis}) $u$ est \deleted{un isomorphisme} bijectif, et \figure{\begin{tikzpicture}[baseline=(t.base)]\node[draw,ellipse,inner sep=4pt] (t) {$u_0$ est un homéomorphisme.};\draw[->] (t.east)--++(0,0.8);\end{tikzpicture}}

(i\deleted{bis}) $u$ est un isomorphisme d'algèbres topologiques
\margin{\figure{\begin{tikzpicture}\draw[->] (0,0)..controls (-0.5,0.5) and (-0.5,0.85)..(0.8,1.1);\end{tikzpicture}}}

(ii) Pour tout $i$, $u_i:B_i\to C_i$ est bijectif, \figure{\begin{tikzpicture}[baseline=(t.base)]\node[draw,ellipse,inner sep=4pt] (t) {$u_0$ est un homéomorphisme};\draw[->] (t.east)--++(0,0.8);\end{tikzpicture}}

(ii bis) pour tout $i$, munissant $B_i$ et $C_i$ des filtrations $I$-adiques, l'homomorphisme $\operatorname{gr}_I(B_i)\to\operatorname{gr}_I(C_i)$ induit par $u\deleted{\unsure{e}}$ est bijectif, \figure{\begin{tikzpicture}[baseline=(t.base)]\node[draw,ellipse,inner sep=4pt,align=left] (t) {$u_0$ un\\homéomorphisme.};\draw[->] (t.east)--++(0,0.9);\end{tikzpicture}}

\margin{On \unsure{suppose} dès le début}

Enfin, lorsque les $\deleted{B}\inserted{C}_i$ sont plats sur les $A_i$, \unsure{les} conditions équivalent également à

(iii) $u_0:B_0\to C_0$ est un \deleted{isomorphisme} \figure{\begin{tikzpicture}[baseline=(t.base)]\node[draw,ellipse,inner sep=4pt] (t) {homéomorphisme.};\draw[->] (t.east)--++(0,0.8);\end{tikzpicture}}

L'équivalence de (i) et (i bis) résulte trivialement du fait que les topologies de $B,C$ sont \deleted{adiques} définies par les puissances d'idéaux contenant $IB$ resp. $IC$. \deleted{\unsure{Celle de}} Evidemment $(\mathrm i\inserted{bis})\Rightarrow(\mathrm {ii})$, pour $(\mathrm {ii})\Rightarrow(\mathrm i\inserted{bis})$ on note que les applications canoniques $B\to\underset{\inserted{\leftarrow}}{\lim}B_i$ et $C\to\underset{\inserted{\leftarrow}}{\lim}C_i$ sont bijectives grâce au fait que les $I_iB$ et $I_iC$ sont fermés (cf rectification à $0_{\inserted{7.2.4.}}$). L'équivalence de (ii) et (ii bis), \deleted{\unsure{qu}} mise pour mémoire, provient du fait que ($I=I_0$ étant un idéal de définition, donc nilpotent dans les $A_i$) les filtrations mises sur $B,C$ sont finies. Enfin, reste à prouver que lorsque \deleted{\unsure{B est p}} $\deleted{B}\inserted{C}_i$ est plat sur $A_i$, alors $u_0$ bijectif implique $u_i$ bijectif ou si on veut $\operatorname{gr}(u)$ bijectif, ce qui est immédiat et sauf erreur figure déjà en lemme quelque part au Chap III (j'espère, avec la généralité qui convient, et bien entendu modulé ...).

\margin{Cf diagramme}
\diagram{\begin{tikzcd}[ampersand replacement=\&]
\operatorname{gr}B\arrow[r,"\mathrm{surj}"] \& \operatorname{gr}C\\
B_0\otimes\operatorname{gr}\unsure{A_i|A'}\arrow[u,"\mathrm{surj}"]\arrow[r,Rightarrow,"\mathrm{bij}"'] \& C_0\otimes\operatorname{gr}\unsure{A_i|A'}\arrow[u,"\mathrm{bij}"']
\end{tikzcd}}

\mspage{9}
